\section[Noun]{Noun (Nennwort)} \label{sec:denoun}
There is a tri-gender system. Each noun can either be male, female or neutral, and has  singular and plural forms that are modified by  four cases. The problem is that there are multiple patterns for the pluralization of a word for a particular gender. I will list these patterns directly in \autoref{tab:denounT1}, cross reference them with the various words you find as you learn. All  nouns  are in case $1$ - \textgerman{\bt{Nominativ}}, which is also their noun root.
\begin{table}[tbh]
    \centering
    \begin{tabular}{|c|c|c|}
       \hline \bt{Männlich} & \bt{Weiblich} & \bt{Sächlich} \\\hline
        Verlag / Verlage & Schachtel / Schachteln & Gehalt / Gehälter \\\hline
        Körper / Körper & Kraft / Kräfte & Essen / Essen \\\hline
        Sturz / Stürze & Mutter / Mütter & Gericht / Gerichte \\\hline
        Rand / Ränder & Bahn / Bahnen & Gebirge / Gebirge \\\hline
        Boden / Böden & Anwältin / Anwältinnen & \textcolor{red}{Auto / Autos} \\\hline
        Gedanke / Gedanken &  &  \\\hline
        \textcolor{red}{Modus / Modi} & & \\\hline
    \end{tabular}
    \caption{Gender based noun Singular / Plural}
    \label{tab:denounT1}
\end{table}

Any German noun must fit into one of the patterns listed in \autoref{tab:denounT1}. For example,  \bt{Gott / Götter} is covered under \bt{Rand / Ränder} as an \textgerman{\bt{\enquote*{Umlaut + er}}} type plural. The red  entries are foreign (loan) words: their list is not comprehensive. Often, they break traditional German rules. Some abstract nouns cannot be pluralized: \bt{Geduld, Schmach}. Some are only plural: \bt{Semesterferien, Leute}. Some nouns also possess multiple plurals or multiple genders, wherein each gender carries a different meaning:
\begin{example}
    \begin{tabular}{c|c|c}
        \bt{Männlich} & \bt{Gehalt} & Content of something in a mixture \\\hline
        \bt{Sächlich} & \bt{Gehalt} & Salary \\\hline
        \bt{Männlich} & \bt{Kiefer} & Jaw \\\hline
        \bt{Weiblich} & \bt{Kiefer} & Pine tree
    \end{tabular}
\end{example}
Sometimes, the multiple genders do not matter and are  related to dialect differences:
\begin{example}
    \begin{tabular}{c|c|c}
    \bt{Wirrwarr} & \bt{Männlich/Sächlich} & Confusion\\\hline
    \bt{Joghurt} & \bt{Männlich/Weiblich/Sächlich} & Curd
    \end{tabular}
\end{example}
In addition, each noun is attributed an article that represents its gender. Imagine a/an/the but with different forms per case and grammatical gender.

There is a special rule with some male nouns called the N-declination, which causes the forms for all cases (except  $1$) in both singular and plural to become the same as the $1$ plural, which always take the end \bt{-en/n}. The single exception to this rule is the word \bt{Herz} which is the only neutral noun to follow it, that too partially. For neutral nouns in particular, there is an interesting class of words that starts with \bt{ge}. These  words  have some special associated meanings: collective marker or the result of an action:
\begin{example}
    \begin{tabular}{c|c|c}
        \bt{Berg} & \bt{Gebirge} & Mountain vs Mountain range \\\hline
        \bt{Wasser} & \bt{Gewässer} & Water vs Water body\\\hline
        \bt{Malen} & \bt{Gemälde} & To paint vs The result of painting (A painting)\\\hline
        \bt{Stein} & \bt{Gestein} & Stone vs Many stones as a group
    \end{tabular}
\end{example}
These nouns normally do not change form, save for due to the \bt{Dativ, Genetiv} case rules. Note that one cannot conclude in general that a noun starting with \bt{ge}  necessarily falls in this category or is neutral in gender, but such exceptions are rare. Ex: \bt{Geschichte (Weib), Gedanke (Männ), Geisel (Männ/Weib)}. 

An interesting thing  about German words is that they are formed simply, and encode the process of formation itself as a meaning, giving them versatility the likes of which has to be seen to be believed. For example: \bt{Anschlag} is formed of \bt{an} meaning `on/to' and \bt{Schlag} meaning `hit'. So naturally, \bt{Anschlag} is the process of hitting (on/towards) something. This can mean, if performed by one or a group of people, an attack or assassination. If the interpretation is something that is hit on something, then a poster or notice. If considered as the effect of the hit on something, then the sound created by  impact. If perceived as the intention to hit towards  something, then the process of getting into shooting position with a shot based weapon. All of these are valid meanings for this word, and most of German works like this. The consequence is that many times, it is possible to understand new words or make new ones that others may understand, even though a machine translator might get stumped.

All nouns, irrespective of their position in a sentence, must be capitalized: \bt{eine Sache, das Gewässer}
\subsection[Case]{Case (Fall)}
There are four cases in German, not including the trivial $8^{th}$, which means that of the seven we considered, the meanings of three have been compressed into one of the others. Also, each gender has a slightly different pattern of form variation. The case equivalence  along with its corresponding rules are listed in \autoref{tab:denounT2}.
\begin{table}[bth]
    \centering
    \begin{tabular}{|c|c|c|}
        \hline \textit{Official name} & \textit{Case codes} & \textit{Noun forms}\\
        \hline\bt{Nominativ} & $1$ & $\vm{R}$ \\\hline
        \bt{Akkusativ} & $2$ & $\vm{R}$ \\\hline
        \bt{Dativ} & $3,4,5,7$ & $\vm{R}(1)$ / $\vm{R}(1)$ + \bt{e}; $\vm{R}(2)$ + \bt{n}\\\hline
        \bt{Genetiv} & $6$ & $\vm{R}(1)$ / $\vm{R}(1)$ + \bt{e(s)}; $\vm{R}(2)$ \\\hline
    \end{tabular}
    \caption{German case equivalence}
    \label{tab:denounT2}
\end{table}

Male and neutral nouns follow the same pattern, and may have an additional \bt{-e} in their \bt{Dativ} singular. For the \bt{Genetiv} singular, they have either \bt{-es/s} added at the end. The only exception to these rules is when N-declination applies: then all word forms except for $\vm{C}(1,1)$ are equivalent to $\vm{R}(2)$. Female nouns do not follow these rules and just use $\vm{R}(1)$ for $\vm{C}(:,1)$. All nouns must have an additional \bt{n} at the end when building the \bt{Dativ} plural. The only exceptions to these rules are loan words that do not really belong, don't worry about them.  \autoref{tab:denounT3} shows  three nouns: one from each gender, with the male displaying N-declination. 
\begin{table}[bth]
    \centering
    \begin{tabular}{|c|c|c|c|c|c|c|}
        \hline $\Diamond$ & \multicolumn{2}{|c|}{Der Rabe} & \multicolumn{2}{|c|}{Die Bö} & \multicolumn{2}{|c|}{Das Land} \\\hline
        \bt{Nom} & Rabe & Raben & Bö & Böen & Land & Länder \\\hline
        \bt{Akk} & Raben & Raben & Bö & Böen & Land & Länder \\\hline
        \bt{Dat} & Raben & Raben & Bö & Böen & Land/Lande & Ländern\\\hline
        \bt{Gen} & Raben & Raben & Bö & Böen & Landes & Länder \\\hline
    \end{tabular}
    \caption{Examples of some case-multiplicity matrices}
    \label{tab:denounT3}
\end{table}
The easiest test to check whether a word follows the N-declination is to check $\vm{C}(4,1)$. If there is an \bt{-en} type ending, that is a guaranteed N-word. The articles \bt{der/die/das} mentioned in \autoref{tab:denounT3} are typically used to mark genders and themselves have declinations.

A tiny disclaimer: Since female nouns have no word alterations in case $6$, often the preposition \bt{von} is used as an indicator with the meaning `of', especially when no article is used with them. Sometimes, this is also done for other nouns when used without articles. In such formations, the noun follows the \bt{Dativ} case as demanded by \bt{von}.
\begin{example}
    Die Rechte von Frauen, Das Haus von Gott, Der König von Bayern...
\end{example}

\subsection[Article]{Article (Geschlechtswort)}
Correspond to the English (a/an)/the,  and are used to mark the gender of a noun. There exist definite, indefinite and negative articles, whose $\vm{G}_{M}(1,1)$ forms are \bt{der, ein, kein} respectively. All of them decline according to  the noun they are attributed to, even if they stand alone in a sentence as their placeholders. All three types follow similar patterns, whose roots are  $\vm{G}(1,:)$. However, it is most convenient to treat the definite article table \ref{tab:denounT4} as a base for deriving  the others. 
\begin{table}[bth]
    \centering
    \begin{tabular}{|c|c|c|c|c|}
        \hline $\Diamond$ & \bt{Männlich} & \bt{Weiblich} & \bt{Sächlich} & \bt{Mehrzahl} \\\hline
        \bt{Nom} $(1)$ & der & die & das & die \\\hline
        \bt{Akk} $(2)$ & den & die & das & die \\\hline
        \bt{Dat} $(3,4,5,7)$ & dem & der & dem & den \\\hline
        \bt{Gen} $(6)$ & des & der & des & der \\\hline
    \end{tabular}
    \caption{Definite article form variation: $\vm{G}$}
    \label{tab:denounT4}
\end{table}

Some examples of article forms: \bt{keine, keinem, einem, einer, eine} and so on. \bt{Kein} and \bt{ein} decline exactly the same, and have an exceptional configuration for $\vm{G}_{M}(1,1)$ and $\vm{G}_{N}(1 \leftrightarrow2, 1)$, where the forms are \bt{ein/kein} rather than \bt{keiner} for example. A structural difference between them is that \bt{ein} is an analogue to a/an and does not possess a plural form. Since this pattern of variation will repeat often, you should have \autoref{tab:denounT4} memorized. Articles can also be used as  $3^{rd}$ person pronouns in certain contexts.

\subsection[Adjective]{Adjective (Eigenschaftswort)}
Adjectives  also double as adverbs. If a non-capitalized word root  ends in any of the following suffixes,  it is an adjective:
\begin{enumerate}
    \item \bt{-ig}: Highly versatile suffix that can be applied on almost any word \begin{example}
        lauig, traurig, dreckig, heutig, eckig\ldots
    \end{example}
    \item \bt{-bar}: Analogue of the `-able' suffix in English and denotes ability. Applied exclusively on verb roots (VR)\begin{example}
        fahrbar, auffindbar, machbar, verwendbar, kletterbar\ldots
    \end{example}
    \item \bt{-lich}: Normally applied on nouns, although  VRs are also allowed. Often have spelling variations\begin{example}
        klanglich, gefährlich, einheitlich, spärlich, sächlich\ldots
    \end{example}
    \item \bt{-isch}: Implies belonging, or the meaning  `in the context of'. These  usually have many germanized words\begin{example}
        akribisch, städtisch, mechanisch, deutsch, theoretisch\ldots
    \end{example}
    \item Pure variants without any particular pattern\begin{example}
        rot, grausam, lau, öde, frisch\ldots
    \end{example}
    \item All \bt{MdV, MdG} forms
\end{enumerate}
This list is not comprehensive but sufficient.

If embedded in a noun unit, adjectives change form according to noun gender, case and multiplicity. The nature of suffixes used mimics \autoref{tab:denounT4}, so even though \autoref{tab:denounT5} may seem daunting, it contains a lot of redundancy.
\begin{table}[bth]
    \centering
    \begin{tabular}{|c|c|c|c|c|}
        \hline \bt{Fall} & \bt{Männlich} & \bt{Weiblich} & \bt{Sächlich} & \bt{Mehrzahl} \\\hline\hline
        \multicolumn{5}{|c|}{\textit{Definite Article} / Bestimmtes Geschlechtswort}\\\hline
        \bt{Nom} & -e & -e & -e & -en \\\hline
        \bt{Akk} & -en & -e & -e & -en \\\hline
        \bt{Dat} & -en & -en & -en & -en \\\hline
        \bt{Gen} & -en & -en & -en & -en \\\hline\hline
        \multicolumn{5}{|c|}{\textit{Possessive pronoun, Indefinite/Negative article}}\\\hline
        \bt{Nom} & -er & -e & -es & -en \\\hline
        \bt{Akk} & -en & -e & -es & -en \\\hline
        \bt{Dat} & -en & -en & -en & -en \\\hline
        \bt{Gen} & -en & -en & -en & -en \\\hline\hline
        \multicolumn{5}{|c|}{\textit{Without article} / Ohne Geschlechtswort}\\\hline
        \bt{Nom} & -er & -e & -es & -e\\\hline
        \bt{Akk} & -en & -e & -es & -e \\\hline
        \bt{Dat} & -em & -er & -em & -en\\\hline
        \bt{Gen} & -en & -er & -en & -er\\\hline
    \end{tabular}
    \caption{Adjective form variation: $\vm{J}$}
    \label{tab:denounT5}
\end{table}

The idea is that whenever an adjective precedes a noun, the nature of use will determine which table is active. For example: \bt{Der mutig(e) Mann} requires $\vm{J}_{M}(1,1)$ in the definite article table on the adjective \bt{mutig}. The same example using the indefinite article: \bt{Ein mutig(er) Mann}. If an adjective stands alone outside a noun unit, none of this is applicable, and only its root  is used. Such sentences always need the existential verb \bt{sein} to work.
\begin{example}
    \begin{tabular}{c|c}
        Dieses Eis ist lecker & Das leckere Eis...\\\hline
        Das Wetter ist kalt & Das kalte Wetter...
    \end{tabular}
\end{example}
An exception occurs when the adjective (root form: comparatives do not follow this) ends in either \bt{l/r}. In such cases, the last and second last letters are reversed before declining: \bt{dunkel/dunkle, übel/üble}. Do note that I am assuming you have enough of a brain to realize that $\vm{J}_{F}(1,1)$  will not end up as \bt{dunklee} but just \bt{dunkle} because an \bt{e} already exists there. As I said earlier, spelling rules are  for ease of pronunciation, so take the tables with a pinch of salt. If it sounds unnatural, there is probably a minor change somewhere.

In some rare cases, some adjectives are  never declined. Ex: \bt{rosa, lila}. This pattern is supposedly found only in colours ending with vowels (although \bt{blau} declines) and words with foreign origin. To be fair, these are so rare, that you can forget this fact and still never have difficulties. Unless of course, colours are a major part of your life.

The general  comparative and superlative root forms in German are formed similar to English and have the same purpose:
\begin{itemize}[leftmargin=25pt]
    \item Comparative: `root' + \bt{-er} (follows the \bt{l/r} principle during formation)
    \item Superlative:  `root' + \bt{-st}
\end{itemize}
So something like \bt{mutig} (courageous) would end up with comparative \bt{mutiger} and superlative \bt{mutigst}. Another example: \bt{dunkel/dunkler/dunkelst}. These higher order adjectives (\bt{Steigerungsformen}) also follow \autoref{tab:denounT5} if kept in the noun unit. To simplify matters, I will mark the associated cases in brackets for the following subsection.

The comparative is used to compare a property of two nouns specifically, or to compare a noun property to its observed average:
\begin{example}
    \begin{tabular}{m{7cm}|m{7cm}}
 \bt{Er ist mutiger($1$) als er} & He is more courageous than him \\\hline
  \bt{Wir wohnen in einer größeren($7$) Stadt} & We live in a \ut{bigger} city
    \end{tabular}
\end{example}
The second sentence means that the city in which we live is bigger than many others (lies on the larger side), but not necessarily bigger than a particular named example.

The superlative expresses extremes:
\begin{example}
    \begin{tabular}{m{7.5cm}|m{7cm}}
  \bt{Er ist am mutigsten($7$)} & He is the most courageous\\\hline
  \bt{Der mutigste($1$) Ritter ist im Gefecht gestorben} & The bravest knight has died in the battle
    \end{tabular}
\end{example}
In the first example, \bt{am} = \bt{an + dem}, and is a short form. These are explained at the very end in \autoref{sec:despco}, and will be explicitly expanded if used.  Another superlative use case would be phrases like `one of the most popular\ldots' In such cases, you use the indefinite case $1$ article for the noun  being graded, and place $\vm{G}(4,2)$ with the superlative directly after it.  Note that the indefinite articles used here are \bt{einer/eine/eines} for the male/female/neutral cases.
\begin{example}
    \begin{tabular}{m{7.5cm}|m{7.5cm}}
        \bt{Einer der stärksten Männer der Welt ist Yujiro Hanma} & One of the strongest men of the world is Yujiro Hanma \\\hline
        \bt{Eine der größten Städte in Deutschland ist München} & One of the biggest cities in Germany is Munich\\\hline
        \bt{Eines der gerechtesten Räder des Weltalls ist das Schicksalsrad} & One of the most just wheels of the universe is the wheel of fate
    \end{tabular}
\end{example}
Nothing would be natural about natural language without exceptional configurations, which appear for single-syllabic adjectives: a syllable roughly being the extent of a word you can pronounce without significantly changing the state (jaw and tongue position) of your mouth. More precisely, it is a distinct speech unit  around a vowel. In such cases, the comparative receives an \bt{Umlaut} on whichever alphabet can receive it: \bt{ä,ö,ü}.
\begin{example}
    \begin{tabular}{c}
    \bt{groß/größer/größt, rot/röter/rötest}
    \end{tabular}
\end{example}
Then are the adjectives beyond rhyme or reason, but English is no better here:
\begin{example}
    \begin{tabular}{c}
    \bt{gut/besser/best, hoch/höher/höchst, gern/lieber/liebst}
    \end{tabular}
\end{example}
However, not all adjectives have comparative and superlative forms, by virtue of their meaning: \bt{leer} (empty), \bt{schwanger} (pregnant), \bt{tot} (dead). These adjectives are the only ones that do NOT double as adverbs.

\subsection[Pronoun]{Pronoun (Fürwort)}
Pronouns  do not need compulsory capitalization within sentences.  We start with the class of personal pronouns. Since German only has two stages of formality: formal/informal, the formal forms have not been explicitly considered in \autoref{tab:denounT6}.
\begin{table}[bth]
    \centering
    \begin{subtable}{\textwidth}
        \centering
    \begin{tabular}{|c|c|c|c|}
        \hline \multicolumn{2}{|c|}{$1_p$} & \multicolumn{2}{|c|}{$2_p$}\\\hline
        \bt{Ein} &\bt{Mehr} & \bt{Ein} & \bt{Mehr}\\\hline
        ich & wir & du & ihr\\\hline
        mich & uns & dich & euch\\\hline
        mir & uns & dir & euch\\\hline
        mein/meine & unser/unsere & dein/deine & euer/eure\\\hline
    \end{tabular}
    \caption{$K = 1_{p}, 2_{p}$}
\end{subtable}
\begin{subtable}{\textwidth}
    \centering
\begin{tabular}{|c|c|c|c|}
        \hline \multicolumn{4}{|c|}{$3_p$}\\\hline
        \bt{Männ} &\bt{Weib} & \bt{Säch} & \bt{Mehr}\\\hline
        er & sie & es & sie \\\hline
        ihn & sie & es & sie\\\hline
        ihm & ihr & ihm & ihnen\\\hline
        sein/seine & ihr/ihre & sein/seine & ihr/ihre\\\hline
    \end{tabular}
    \caption{$K = 3_{p}$}
\end{subtable}
\caption{Pronoun form variation: $\vm{P}[a,In,K]$}
\label{tab:denounT6}
\end{table}

For  the polite second person, singular and plural for any particular case are the same. Use the plural forms of the given third person, capitalize the first letter, and there you have the polite second person. Important is that in the polite variant, the pronoun MUST be capitalized irrespective of its place in the sentence. This matrix will be stacked in the third dimension above $\vm{P}[a,In,2_p](:,:,0)$ as $\vm{P}[a,In,2_p](:,:,\mathbb{H})$. Case $6$ personal pronouns take on the gender of the noun  being possessed:
\begin{example}
    \begin{tabular}{c|c}
        \bt{Diese Luft ist ihre} & This air is hers\\\hline
        \bt{Die Fahrräder sind ihre} & The bicycles are theirs\\\hline
        \bt{Das Buch ist mein} & The book is mine
    \end{tabular}
\end{example}
The forms for the `personal' and `possessive' pronoun classes are not distinguished in \bt{Deutsch}, but are clearly different. The roots for $\vm{P}[a,De,K\{M\}]$ are based upon the corresponding form in $\vm{P}[a,In,K](4,M)$, where $K$ represents the person being considered and $M$ its multiplicity. These roots  decline according to the negative article \bt{kein}. The polite variants $\vm{P}[a,De,2_p\{1\leftrightarrow\mathbb{M}, \mathbb{H}\}]$ require compulsory capitalization. An example for $\vm{P}[a,De,2_p\{\mathbb{M},0\}]$ is in \autoref{tab:denounT7}.
\begin{table}[bth]
    \centering
    \begin{tabular}{|c|c|c|c|}
        \hline\bt{Männ} & \bt{Weib} & \bt{Säch} & \bt{Mehr}\\\hline
        euer & eure & euer & eure\\\hline
        euren & eure & euer & eure\\\hline
        eurem & eurer & eurem & euren\\\hline
        eures & eurer & eures & eurer\\\hline
    \end{tabular}
    \caption{$2_p$ Plural possessive pronoun}
    \label{tab:denounT7}
\end{table}
Some clarifying examples:
\begin{example}
    \begin{tabular}{c|c}
        \bt{(Meine) Fahrräder haben keine Makel} & My bicycles have no faults\\\hline
        \bt{Sie hat sie unter (ihrer) Obhut} & She has her under her care\\\hline
        \bt{Alle lieben (Ihre) Vorträge} & All love his/her lectures
    \end{tabular}
\end{example}
Reflexive pronouns are encompassed under \bt{sich} and its forms as listed in \autoref{tab:denounT8}. These pronouns just imply self-referencing, and cannot be used with every verb. Their placement is always exactly after the conjugating verb or its illusion (in case it is exiled to the end of the sentence).
\begin{table}[bth]
    \centering
    \begin{tabular}{|c|c|c|c|c|c|c|}
        \hline $\Diamond$ & \multicolumn{2}{|c|}{$K=1_p$} & \multicolumn{2}{|c|}{$K=2_p$} & \multicolumn{2}{|c|}{$K=3_p$}\\\hline
        \bt{Akk} & mich & uns & dich & euch & sich & sich \\\hline
        \bt{Dat} & mir & uns & dir & euch & sich & sich \\\hline
    \end{tabular}
    \caption{Reflexive pronouns: $\vm{P}[r,De,K]$}
    \label{tab:denounT8}
\end{table}
\begin{example}
    \begin{tabular}{m{7cm}|m{7cm}}
        %\bt{Er stellt sich hinter mich} & He positioned himself behind me\\\hline
        \bt{Ich wasche mir die Hände}& I wash \ut{using myself the hands}\\\hline
        \bt{Wir treffen uns am Montag} & We \ut{are meeting us} on Monday\\\hline
        \bt{Ich bedanke mich fürs\{für + das\} Zuschauen} & I \ut{give the thanks within myself to} \ut{(those who saw)} for the viewing
    \end{tabular}
\end{example}
The last example  uses a special  construction used for fixed reflexiveness, something not found in English. It essentially means, `the thanks came from him' but it never mentions to whom it is given, implying it as obvious information. Here, the choice of the reflexive pronoun must additionally match  the fixed case code for the concerned usage. In this case, \bt{Ich bedanke mich} demands case $2$, so \bt{ich bedanke mir} is invalid. Occurrences are rare and the advice is to memorize these. Formal \bt{Sie} reflexives are  $\vm{P}[r,De,3_p]$ forms: these do not need capitalization.

Indefinite pronouns change forms according to corresponding article templates in both dependent and independent uses, but these changes are not consistently assigned. Some indefinite pronouns  can be used exclusively for objects, some exclusively for people, while most others fit for both. \autoref{tab:denounT9} documents various examples along with their followed  form variations. If any adjectives are used in the noun unit during such dependent use, they can be  declined in one of two ways:
\begin{enumerate}
    \item Assume the pronoun is the template word, and decline normally
    \item Use `without article' declinations
\end{enumerate}
The first rule merely assumes that the pronoun is structurally the same as the word, whose form template it follows. Both these rules are mutually exclusive, but have been assigned rather randomly. To mark pronouns which follow the second rule for adjective form variation, an asterisk (*) is added to their entries in \autoref{tab:denounT9}.
\begin{example}
    \begin{tabular}{c|c}
        [Aus (welchem) seltsamen Grund] benutzt du deine Sphäre nicht? & der\\\hline
        [Für (welche) seelische Flamme] möchtest du dich zerkleinern? & die\\\hline
        [(Manche) teuere Flaschen] sind es nicht wert & die \\\hline
        [(Solche) einzelne Menschen] sind meist Feiglinge & der
    \end{tabular}
\end{example}
The noun units and corresponding indefinite pronouns have been marked. It is also important to note that some of these pronouns may also directly be used as normal adjectives: \bt{Eine solche Sache, die wenigen Leute}\ldots clearly, the separation is fluid, and the only practically important thing is the applicable structure per use case.
\begin{table}[bth]
    \centering
    \begin{subtable}{\textwidth}
        \centering
        \begin{tabular}{|c|c|c|}
            \hline $\diamondsuit$ & \bt{Dependent} & \bt{Independent}\\\hline
            Alle & die (Mehrzahl)& die (Mehrzahl)\\\hline
            Einige* & die (Mehrzahl) & die (Mehrzahl)\\\hline
            Irgendeiner & ein/eine & der/die/das (Einzahl)\\\hline
            Jeder & der/die/das (Einzahl) & der/die/das (Einzahl)\\\hline
            Keiner & kein/keine & der/die/das \\\hline
            Mancher* & der/die/das & der/die/das\\\hline
            Solcher & der/die/das & der/die/das \\\hline
            Derartig* & Ohne Gesch.wort in \autoref{tab:denounT5} & Ohne Gesch.wort in \autoref{tab:denounT5}\\\hline
            Mehrere* & die (Mehrzahl) & die (Mehrzahl)\\\hline
        \end{tabular}
        \caption{Universally applicable}
    \end{subtable}
    \begin{subtable}{\textwidth}
        \centering
        \begin{tabular}{|c|c|c|}
        \hline $\clubsuit$ & \bt{Dependent} & \bt{Independent}\\\hline
        Nichts & ! (Ungültig) & - (Keine Formänderung)\\\hline
        Alles & ! & das (Einzahl)\\\hline
        Welcher & der/die/das & der/die/das \\\hline
        Irgendwelcher & der/die/das & der/die/das\\\hline
        Viel/Viele* & - & - \\\hline
        Wenig & ! & - \\\hline
        Etwas & ! & -\\\hline
        \end{tabular}
        \caption{Only inanimate objects}
    \end{subtable}
    \begin{subtable}{\textwidth}
        \centering
        \begin{tabular}{|c|c|c|}
            \hline $\spadesuit$ & \bt{Dependent} & \bt{Independent}\\\hline
            Man & ! & ein (Männ, kein Fall $6$)\\\hline
            Jemand & ! & ein (Männ, kein Fall $6$)\\\hline
            Irgendjemand & ! & ein (Männ, kein Fall $6$)\\\hline
            Niemand & ! & ein (Männ, kein Fall $6$)\\\hline
            Irgendwer & ! & der ($\vm{Q}_{M}(:,1)$) \\\hline 
        \end{tabular}
        \caption{Only people}
    \end{subtable}
    \caption{Indefinite pronouns: $\vm{P}[b,De/In]$}
    \label{tab:denounT9}
\end{table}

Some supplementary explanations to \autoref{tab:denounT9} will help digest it better:
\begin{enumerate}
    \item \bt{Alles/Alle} - \bt{Alles} represents a whole of something, and strictly refers to non-living entities, whereas \bt{Alle} has no such restriction and represents multiple wholes. \begin{example}
        \begin{tabular}{c|c}
            \bt{Alles war zerstört} & Everything was destroyed \\\hline
            \bt{Alle waren nass} & All were \ut{being} wet 
        \end{tabular}
    \end{example}
    \item \bt{Derartig} is an abbreviation of \textgerman{\enquote{\bt{von dieser Art}}} and declines like an adjective irrespective of whether it is used as a pronoun or adjective. The initial \bt{der} is not to be touched, only the endings vary.
    \item The prefix \bt{irgend-} (some) may be added to most of the W-words like `when, why, who\ldots' to create `someone, something, somebody\ldots' and so on. Using \bt{nirgend-} negates the corresponding word: `nobody, nowhere\ldots' \bt{Irgend-/Nirgend-} are merely  significance markers: using them on any word only decreases its perceived significance. Ex: \begin{example}
        \begin{tabular}{c|m{7cm}}
            \bt{Jemand ist angekommen} & Someone (Unknown) has arrived \\\hline
            \bt{Irgendjemand ist angekommen} & Someone (Unknown and irrelevant) has arrived
        \end{tabular}
    \end{example}
    These \bt{irgend-} compounds never decline. The only ones who do are documented in \autoref{tab:denounT9}. Commonly appended pronouns are: \bt{etwas, wo, wie, was, wann}.
    \item The diligent reader must have realized that the categories of `pronoun, adjective, article' are not completely independent. As long as the meaning supports it, any word may be a member of multiple  classes, where it will follow the respective structure. Ex: \begin{example}
        \begin{tabular}{c|c|c|c}
            \bt{Base} & \bt{Pronoun} & \bt{Adjective} & \bt{Article}\\\hline
            ein & einer & ein & ein \\\hline
            viel & viel & viele & ! 
        \end{tabular}
    \end{example}
    \bt{Viel/Viele} are the same as  much/many, and only mark countability. Since no emphasis is marked, article usage becomes invalid.
    \item \bt{Wer} requires hook word forms which will be introduced in \autoref{sec:destce}.
\end{enumerate}

\autoref{tab:denounT10} details the  demonstrative pronouns along with their corresponding form variations.
\begin{table}[bth]
    \centering
    \begin{tabular}{|c|c|c|c|}
        \hline $\heartsuit$ & \bt{Dependent} & \bt{Independent} & \ut{Meaning}\\\hline
        Dies & der/die/das & kein/keine & \textit{This} \\\hline
        Das & ! & - & \textit{This/That}\\\hline
        \textcolor{magenta}{Derselbe} & \autoref{tab:denounT5}, \autoref{tab:denounT4} & \autoref{tab:denounT5}, \autoref{tab:denounT4} & \textit{the same}\\\hline
    \end{tabular}
    \caption{Demonstrative pronouns: $\vm{P}[c,De/In]$}
    \label{tab:denounT10}
\end{table}
Both \bt{das} and \bt{dies} have the same meaning, but with different emphases. \bt{Das} is rather neutral and general, while \bt{dies} is formal and much more concrete. \textcolor{magenta}{Derselbe} is a fusion of the definite article and the adjective \bt{selb}, meaning `same'. This word follows all noun unit rules for articles and adjectives, and implies that something is exactly something else, which should not be confused with \bt{das gleiche}, which implies that something is exactly equal to something else. \bt{Derselbe} implies that the two entities being considered are one and the same, not merely equal.
\begin{example}
    \begin{tabular}{m{7.5cm}|m{7.5cm}}
        \bt{Grober Kies liegt auf diesem Strand} & Coarse gravel lies on this beach\\\hline
        \bt{Das kann unseren Einsatz gefährden} & This could \ut{our mission endanger}\\\hline
        \bt{Zwillinge entstehen aus demselben Mutterleib, und genießen oft die gleiche Gunst} & Twins emerge out of the same womb, and \ut{enjoy often} the same favour
    \end{tabular}
\end{example}
An exact analogue to a far-away marker like `that' does not exist.  Adverbs like \bt{drüben, dort, da}  concretely mark distance, and are used with \bt{das} to get `that'. Ex: \bt{das dort, dort, drüben, das da, da drüben\ldots}

\subsection[Joining]{Joining (Wortverkettung)}
I hinted earlier that \bt{Deutsch} uses a lot of concatenated words that encode meanings via processes. Compared to the sheer amount of such compounds, the rules concerning them are surprisingly compact:
\begin{enumerate}
    \item Add \bt{-s} to the end of the first word and join
    \item Add \bt{-en/-n} to the end of the first word and join
    \item If the first word is a verb, eliminate its infinitive \bt{-en} and join
    \item Join
\end{enumerate}
These four rules are mutually exclusive (obviously), but there is no real pattern as to which applies where. A lot of \bt{Deutsch} has old relics still in use,  even though the language evolved and nobody uses the old rules anymore. New words are formed using new rules, and the ones which used the old rules remain as exceptions\ldots of which there are many. Giving some examples is all I can do for this section:
\begin{example}
    \begin{tabular}{c|c|c}
        Anmeldung(s)/zeit/raum & Flug/zeug & Gär(-en)/kessel\\\hline
        Tag(es)/licht & Schnecke(n)/haus & Wiedergabe/geschwindigkeit\\\hline
        Krank(en)/wagen & Tür/störung & Fahr(-en)/strecke
    \end{tabular}
\end{example}
You will always join the singular of the first word. The compound can then be treated like a regular noun. This compound noun takes the gender and word forms of the last noun in the chain. Oh, and there is technically no limit to how many of those you can smash together as long as it makes sense. This convenience is above the mental capacities of some people, which is why German gets insulted by the inept for its long words, complicated nature and aggressiveness. It is also likely because of these long words that Germans are habituated to speaking fast, else sentences would take a while to complete.

Mostly, the first word has cases $6$ or $4$ if the compound were to be expanded. Keep in mind  that the associated case is not necessarily unique, since the expansion is not unique. It just has to make sense. This feature `cleans' the language by introducing roughness, which is why English speakers  struggle initially.
\begin{example}
    \begin{tabular}{m{7cm}|m{7cm}}
        \bt{Anmeldungszeitraum}  & Registration time period\\\hline
        \bt{Flugzeug} &  Airplane (Flight thing)\\\hline
        \bt{Wiedergabegeschwindigkeit} &  Playback speed
    \end{tabular}
\end{example}
If the meaning allows it, adjectives/nouns  can also be pre-joined to other adjectives. In such cases, \ut{only} joining rule `4' is allowed:
\begin{example}
    nagelneu, blitzsauber, uralt, bildschön, allerschönst, größtmöglich\ldots
\end{example}
\subsection[Numbers]{Numbers (Zahlen)}
Before starting this subsection, it must be declared that the gender of German numbers, if at all required is always \bt{sächlich}. Let us begin with pure numbers. These  never show form variation, even when placed in noun units (except \bt{eins}, which  becomes the indefinite article \bt{ein}). Such numbers are used to portray time. If written as digits, \bt{Deutsch} uses a comma as a decimal point.
\begin{example}
eins, drei, zwanzig, zweiundzwanzig, achtundvierzig, einundfünfzig, elf\ldots
\end{example}
To count ranks, there are two rules:
\begin{enumerate}
    \item $\{1,2,\ldots, 19\}$: Add \bt{-t} at the end of the pure number. $1$ and $3$ are the only exceptions with counters \bt{erst,  dritt} respectively.
    \item $\{20,21,\ldots\}$: Add \bt{-st} at the end of the pure number.
\end{enumerate}
These rank counter root forms behave like adjectives, and are subject to their rules. It is possible to say things like `as a couple, as a quintuple' by pairing \bt{zu} and the respective rank root. This usage does not undergo form variation.
\begin{example}
    zu zweit, zu dritt, zu fünft, zu sechst\ldots
\end{example}

To define  multiplicities,  the suffix \bt{-fach} (times) is added after  the factor $f$ as defined in \autoref{eq:factor}. $[f]$ is represented by pure numbers. Fractional bases use the suffix \bt{-el} at the end of the corresponding rank counter, with their multiplicity (as a pure number)  appended at the beginning to make $\{f\}$. $f$ is created by joining both of these  in the order $([f], \{f\})$ with a  dash between the components to ensure readability:
\begin{example}
    \begin{tabular}{c|c}
    \bt{Integral} & (einzeln), zweifach, dreifach, vierfach, einundfünfzigfach\ldots\\\hline
    \bt{Fractional} & (einhalb \{$\frac{1}{2}$\}), dreiviertel ($\frac{3}{4}$), neundreiundachtzigstel ($\frac{9}{83}$)\ldots\\\hline
    \bt{Mixed} & zwei-eindrittel-fach ($2\frac{1}{3}$), fünfundfünfzig-dreiviertel-fach ($55\frac{3}{4}$)\ldots
    \end{tabular}
\end{example}
All of these behave like adjectives. `Single' can be written as \bt{einfach} but is usually represented by \bt{einzeln}, while \bt{einfach} means `simple'. 

Verbs that represent the action of multiplying by a factor have the pattern:
\begin{syntax}
    \begin{tabular}{c|c}
    Multiplication: &\bt{ver +} Factor \bt{+ fach + en}\\\hline
    Division: & Pure number \bt{+ teilen}
    \end{tabular}
\end{syntax}
This then gives rise to verbs like: \bt{verzweifachen/verdoppeln,  verzweiundzwanzigfachen, verzwei-eindrittel-fachen, vervierelftel-fachen, (halbieren)\ldots}. Be mindful that fractional components always need \textit{CoN} separation (-). Pure number division verbs denote the division of an object into its argument number of parts: \bt{zweiteilen, vierteilen, zweiundzwanzigteilen\ldots}. These verb classes are regular and use the helping verb \bt{haben}.